\chapter{Proposta}\label{chp:PROPOSTA}

Este capítulo apresenta o problema central enfrentado na construção de um chatbot de busca sobre as atas institucionais da UFRRJ, com a geração de consultas adequadas ao motor de busca a partir da interação em linguagem natural com o usuário  e a solução proposta para esse problema, baseada em um roteiro guiado de diálogo que traduz as respostas do usuário em parâmetros estruturados de uma consulta ao Elasticsearch.

\section{Problema: Gerando Consultas Usando Chatbots}

\subsection{Delimitação da Abordagem: Exclusão de Outros Tipos de Chatbot}

Como discutido na Seção~\ref{sec:comparativo-chatbots} do Capítulo~\ref{chp:FUNDAMENTACAO}, chatbots impulsionados por IA, de IA generativa ou híbridos têm maior capacidade de interpretar perguntas livres do que um chatbot baseado em regras. A Tabela~\ref{tab:comparativo-chatbots} (Seção~\ref{sec:comparativo-chatbots}) resume essa comparação à luz dos requisitos do projeto, e a justificativa teórica para a escolha de um chatbot baseado em regras já foi apresentada naquele capítulo. Cabe aqui retomar essa decisão sob o ângulo específico do problema de geração de consultas: o risco prático de uma resposta não embasada no conteúdo real das atas.

No sistema proposto, a resposta apresentada ao usuário deve corresponder a um documento institucional efetivamente existente: a ata recuperada do índice. Um chatbot de IA generativa, ao gerar texto livremente a partir de um modelo de linguagem, pode produzir uma resposta gramaticalmente correta e aparentemente segura, mas sem garantia de que seu conteúdo corresponda ao que está de fato registrado nas atas. Esse tipo de erro, conhecido como alucinação, é particularmente grave neste domínio: o usuário poderia receber uma informação inexistente ou incorreta sobre uma decisão acadêmica real, como a data de uma reunião, o resultado de uma progressão funcional ou o teor de um processo.

Esse risco já se concretizou fora do contexto acadêmico, em uma disputa judicial real. Em fevereiro de 2024, o \textit{British Columbia Civil Resolution Tribunal}, no Canadá, responsabilizou a empresa aérea Air Canada porque seu chatbot informou incorretamente a um passageiro que ele poderia solicitar, de forma retroativa, uma tarifa reduzida por luto familiar — política que a empresa, na realidade, não permitia. O tribunal rejeitou o argumento de que o chatbot seria uma entidade distinta da empresa e determinou que a Air Canada honrasse a informação incorreta fornecida por seu próprio sistema \cite{aircanada-cbc, aircanada-aba}. O caso ilustra, de forma concreta, que uma resposta tecnicamente convincente, mas não embasada na fonte de verdade correta, tem consequências reais para quem a recebe e para quem a fornece.

Esse argumento reforça, no contexto deste trabalho, a opção por um chatbot baseado em regras combinado a uma busca textual determinística sobre o Elasticsearch: o sistema proposto só pode retornar como resposta o que de fato está indexado a partir do conteúdo das atas, sem gerar texto novo. Os capítulos seguintes detalham como essa restrição é traduzida em uma consulta estruturada (Seção~\ref{sec:problema-gerar-consulta}, a seguir) e implementada na integração com o Elasticsearch (Capítulo~\ref{chp:EXPERIMENTOS}).

\subsection{O Papel da Consulta no Processo de Busca}

Como discutido na Seção~\ref{sec:problema-busca} do Capítulo~\ref{chp:FUNDAMENTACAO}, um sistema de Recuperação da Informação depende de duas etapas: a indexação prévia da coleção de documentos e a busca, que utiliza uma consulta para localizar e ranquear os documentos relevantes dentro do índice \cite{manning2008ir}. A consulta é, portanto, a interface entre a necessidade de informação do usuário e o índice construído sobre a coleção: é a partir dela que o motor de busca seleciona quais campos consultar, com quais operadores e com qual peso cada termo deve contribuir para a pontuação final.

No sistema proposto, essa consulta não é digitada diretamente pelo usuário em uma caixa de busca livre. Em vez disso, ela é construída a partir das respostas fornecidas ao chatbot ao longo de um diálogo guiado e, somente então, repassada ao componente responsável por montar a consulta estruturada para o Elasticsearch. Essa separação entre a interação conversacional e a montagem da consulta é o ponto central da proposta: cabe ao chatbot decidir quais informações coletar e em que formato, e cabe ao motor de busca interpretar essas informações como critérios de relevância. Como cada critério respondido pelo usuário se traduz em um parâmetro adicional da consulta, quanto mais critérios informados, mais restritiva a busca, aumentando as chances de que a ata retornada seja, de fato, a que o usuário procura.

\subsection{O Problema de Gerar a Consulta sem Conhecer o Índice}\label{sec:problema-gerar-consulta}

Um chatbot baseado em regras conduz a conversa por um fluxo pré-definido, mas o conteúdo das mensagens trocadas com o usuário é, em geral, texto livre: o usuário pode digitar o nome de um professor com ou sem acentos, abreviar o nome de um departamento, descrever um período de forma textual (``de março a abril de 2026'') ou usar termos coloquiais para se referir a um assunto. Esse vocabulário não tem relação direta com a forma como o conteúdo das atas está organizado no índice do Elasticsearch, que distingue, por exemplo, o texto integral de cada documento de metadados estruturados como a data a que a ata se refere \cite{IBMElasticsearch}.

Alguns exemplos concretos ilustram esse descompasso. Ao ser perguntado sobre o período da reunião, o usuário pode responder ``de março a abril de 2026'' em vez de informar datas em um formato estruturado; essa resposta precisa ser interpretada como um intervalo e direcionada a um filtro de data, e não ao conteúdo textual da busca. Ao ser perguntado sobre o docente envolvido, o usuário pode digitar apenas ``professor que apresentou sobre TCC'', sem citar um nome próprio, misturando no mesmo campo um termo descritivo de assunto com o critério de docente. Já uma resposta como ``reunião sobre progressão docente'' mistura, em uma única frase, um termo relacionado ao critério de progressão funcional com uma referência genérica ao tipo de evento, sem indicar se o termo relevante para a busca é ``progressão'' ou a frase completa. Em todos esses casos, o texto livre fornecido pelo usuário precisa ser mapeado para o campo correto do índice — conteúdo textual ou filtro estruturado de data — e associado ao peso de relevância adequado ao critério correspondente, e não tratado de forma uniforme.

Essa desconexão entre o vocabulário livre do usuário e a organização do índice é o problema central enfrentado na geração da consulta: se cada resposta do usuário fosse simplesmente concatenada e enviada como uma única \textit{string} de busca, o motor de busca trataria todos os termos com o mesmo peso e buscaria apenas no conteúdo textual das atas, perdendo a oportunidade de usar informações estruturadas para filtrar resultados, e de atribuir pesos diferentes a termos que, semanticamente, têm relevância distinta (por exemplo, o nome de um discente normalmente é mais discriminativo do que uma palavra de um assunto genérico).

\subsection{Definindo o que Pode ou Não Ser Buscado}

A solução para essa desconexão começa pela delimitação explícita do que pode ser buscado. O fluxo principal do chatbot conduz o usuário por um conjunto fechado de critérios de busca, cada um correspondendo a um aspecto recorrente do conteúdo das atas: o departamento responsável pela reunião, a disciplina mencionada, os docentes e discentes envolvidos no assunto tratado, o período em que a ata foi produzida, o número de um processo ou edital referenciado e termos relacionados a progressões funcionais ou acadêmicas. Além desses critérios, o roteiro reserva um campo de assunto geral, utilizado tanto para descrições livres do tema da ata quanto como destino de qualquer informação que o usuário forneça e que não se enquadre nos demais critérios.

Cada um desses critérios corresponde, na consulta enviada ao Elasticsearch, a um componente de busca sobre o conteúdo textual das atas, com um tipo de correspondência e um peso de relevância próprios: critérios mais específicos e discriminativos recebem maior peso e exigem correspondência mais estrita, enquanto critérios mais genéricos recebem peso menor e admitem maior tolerância a variações textuais. O período informado pelo usuário, por sua vez, não contribui para a pontuação de relevância: ele é tratado como um filtro sobre a data da ata, restringindo o conjunto de documentos candidatos sem afetar o \textit{score} atribuído pelo BM25. Dessa forma, delimita-se de forma explícita o que pode ser buscado e como cada resposta do usuário se traduz em um componente da consulta ao Elasticsearch, fechando a lacuna entre o vocabulário livre da conversa e a organização do índice. O mapeamento detalhado de cada critério para os parâmetros de implementação correspondentes é apresentado no Capítulo~\ref{chp:EXPERIMENTOS} (Seção~\ref{sec:mapeamento-criterios}).

\subsection{Solução Proposta: Roteiro Guiado de Diálogo}

A solução proposta consiste em um roteiro guiado de diálogo que conduz o usuário por perguntas relativas a cada um dos critérios de busca apresentados na seção anterior. Cada resposta do usuário é armazenada ao longo da conversa, sem que ele precise conhecer a forma como o conteúdo das atas está organizado no índice ou a sintaxe de consulta do Elasticsearch.

Ao final do roteiro, as respostas coletadas passam por uma etapa de consolidação, na qual a resposta sobre o departamento é separada em nome e sigla,  permitindo que a busca encontre a ata tanto pelo nome por extenso quanto pela forma abreviada ,e a resposta sobre o período é interpretada como um intervalo de datas; quando o texto fornecido não pode ser interpretado como uma data, ele é redirecionado para o critério de assunto geral, evitando que a informação seja descartada. As respostas consolidadas são então traduzidas em uma consulta estruturada, na qual cada critério informado se torna um componente de busca com seu próprio peso de relevância, e o período, quando informado, se torna um filtro sobre a data da ata. Essa consulta é enviada ao motor de busca, que retorna os documentos ranqueados por BM25 \cite{IBMElasticsearch}. Por fim, sobre os documentos retornados é aplicado um cálculo adicional de proporcionalidade com a fração dos termos informados pelo usuário que efetivamente aparecem no título e no trecho de cada ata, usado para reordenar os melhores resultados antes de apresentá-los.

Esse roteiro guiado torna o processo de geração da consulta transparente em dois sentidos: para o usuário, que é conduzido por perguntas em linguagem natural sem precisar conhecer a organização do índice; e para o desenvolvedor, que pode rastrear de forma direta como cada resposta do usuário se transforma em um componente da consulta enviada ao Elasticsearch, facilitando a manutenção e a extensão do sistema e tornando a adição de um novo critério de busca simples, exigindo apenas a inclusão de uma nova etapa no roteiro de diálogo e do componente de busca correspondente na consulta. A implementação concreta desse roteiro, incluindo as ferramentas utilizadas e os detalhes da consulta ao Elasticsearch, é descrita no Capítulo~\ref{chp:EXPERIMENTOS}.
